\chapter*{Danh mục các từ viết tắt, khái niệm và thuật ngữ}
\addcontentsline{toc}{chapter}{Danh mục các từ viết tắt, khái niệm và thuật ngữ}

\renewcommand{\arraystretch}{1.35}
\begin{longtable}{@{}p{2.6cm}p{5.4cm}p{6.2cm}@{}}
\toprule
\textbf{Viết tắt} & \textbf{Từ đầy đủ} & \textbf{Tiếng Việt / giải thích} \\
\midrule
\endfirsthead
\toprule
\textbf{Viết tắt} & \textbf{Từ đầy đủ} & \textbf{Tiếng Việt / giải thích} \\
\midrule
\endhead

GUI & Graphical User Interface & Giao diện đồ hoạ người dùng \\
VLM & Vision-Language Model & Mô hình thị giác và ngôn ngữ \\
LLM & Large Language Model & Mô hình ngôn ngữ lớn \\
SFT & Supervised Fine-Tuning & Tinh chỉnh có giám sát \\
LoRA & Low-Rank Adaptation & Tinh chỉnh hạng thấp \\
QLoRA & Quantized LoRA & LoRA trên trọng số đã lượng tử hoá \\
NF4 & 4-bit NormalFloat & Định dạng lượng tử hoá 4 bit \\
AC & AndroidControl & Bộ dữ liệu điều khiển Android (NeurIPS 2024) \\
AITW & AndroidInTheWild & Bộ dữ liệu Android quy mô lớn (NeurIPS 2023) \\
OCR & Optical Character Recognition & Nhận dạng chữ trong ảnh \\
VH & View Hierarchy & Cây phân cấp khung nhìn của màn hình Android \\
a11y & Accessibility (tree) & Cây trợ năng; nguồn khai báo vai trò và tên phần tử \\
REG & Referring Expression Generation & Sinh biểu thức quy chiếu (mô tả phân biệt một đối tượng) \\
MDE & Minimum Detectable Effect & Mức chênh lệch nhỏ nhất còn phân biệt được với nhiễu \\
KTC & (Confidence Interval) & Khoảng tin cậy \\
dp & density-independent pixel & Đơn vị độ dài độc lập mật độ điểm ảnh trên Android \\
pp & percentage point & Điểm phần trăm (chênh lệch tuyệt đối giữa hai tỉ lệ) \\
ORPO & Odds Ratio Preference Optimization & Mục tiêu ưu tiên gắn thẳng vào hàm mất mát có
giám sát, không cần mô hình tham chiếu \\
\midrule
\multicolumn{3}{@{}l@{}}{\textit{Thuật ngữ riêng của luận văn}} \\
\midrule
\multicolumn{2}{@{}p{8.2cm}}{\textit{executability} (tính thực thi được)}
& Thước đo mà luận văn đề xuất. Một câu hướng dẫn được tính là đúng nếu đưa câu đó cho
một mô hình định vị độc lập thì nó chỉ trúng phần tử mà người dùng đã chạm \\
\multicolumn{2}{@{}p{8.2cm}}{Mô hình định vị phần tử (\textit{GUI grounding model})}
& Mô hình nhận một câu mô tả cùng một ảnh màn hình và trả về toạ độ của phần tử được mô
tả. Ở đây nó đóng vai \textbf{dụng cụ đo}, không phải một hệ thống đem ra so sánh \\
\multicolumn{2}{@{}p{8.2cm}}{Dòng khai báo (\textit{descriptor})}
& Một dòng có cấu trúc gồm bốn ô, \code{[vai trò | tên | toạ độ | dấu hiệu phân biệt]},
mà mô hình phải viết ra trước khi viết câu hướng dẫn. Dòng này bị cắt bỏ trước khi chấm \\
\multicolumn{2}{@{}p{8.2cm}}{Ô Voronoi}
& Vùng trên màn hình gồm những điểm nằm gần phần tử cần chạm hơn bất kỳ phần tử nào khác.
Luận văn dùng vùng này để xác định thế nào là trỏ trúng, thay cho một ngưỡng sai số cố
định \\
\multicolumn{2}{@{}p{8.2cm}}{Bước chạm}
& Bước thao tác mà người dùng chạm vào màn hình, tức bước có toạ độ đích. Toàn bộ điểm số
của luận văn tính trên nhóm bước này \\
\multicolumn{2}{@{}p{8.2cm}}{Phần tử cần chạm}
& Nút, ô nhập hoặc mục danh sách mà người dùng đã chạm ở bước đang xét, tức thứ mà câu
hướng dẫn phải chỉ ra được \\
\multicolumn{2}{@{}p{8.2cm}}{Quần thể chấm}
& Tập $4.463$ bước chạm của tập kiểm, dùng chung cho mọi nhánh, nên hai nhánh bất kỳ luôn
so được với nhau trên cùng những bước \\
\multicolumn{2}{@{}p{8.2cm}}{Trần và sàn của thước}
& Trần là điểm mà câu do người viết đạt được, tức mức cao nhất mà thước còn cho. Sàn là
điểm mà một câu không mang thông tin gì vẫn đạt được. Mọi điểm số đọc trong dải giữa hai
mức đó \\
\multicolumn{2}{@{}p{8.2cm}}{Nhánh nền và nhánh xử lý}
& Nhánh nền là nhánh không có thành phần đề xuất, dùng làm mốc so. Nhánh xử lý là nhánh
có thành phần đó. Chênh lệch giữa hai nhánh chính là đại lượng cần đo \\
\multicolumn{2}{@{}p{8.2cm}}{Hạt giống}
& Số dùng để khởi tạo phần ngẫu nhiên của một lượt huấn luyện. Hai lượt chỉ khác hạt
giống là hai lần dựng lại cùng một hệ thống, nên chênh lệch giữa chúng là nhiễu \\
\multicolumn{2}{@{}p{8.2cm}}{Điểm lưu (\textit{checkpoint})}
& Bản trọng số mô hình được ghi ra đĩa tại một thời điểm của lượt huấn luyện \\
\multicolumn{2}{@{}p{8.2cm}}{Lát cắt}
& Một cách chia quần thể chấm thành các nhóm nhỏ để đọc kết quả riêng trên từng nhóm, ví
dụ chia theo độ dài câu hoặc theo số phần tử trên màn hình \\
\multicolumn{2}{@{}p{8.2cm}}{Phép so đối chiếu}
& Một chênh lệch đã biết chắc dấu và đủ lớn, dùng để kiểm xem một cách đo mới có còn nhìn
thấy những khác biệt thật hay không \\
\multicolumn{2}{@{}p{8.2cm}}{MIN-DESC}
& Chặng huấn luyện thứ hai, nối tiếp nhánh S2. Mô hình học từ những cặp câu chỉ khác nhau
ở dòng khai báo, còn phần câu hướng dẫn thì giống hệt nhau (Mục~\ref{sec:mindescthietke}) \\
\multicolumn{2}{@{}p{8.2cm}}{CE2-S2}
& Nhánh đối chứng của MIN-DESC, dùng để tách xem phần cải thiện nào là của mục tiêu huấn
luyện mới và phần nào chỉ do được học thêm \\
\bottomrule
\end{longtable}
\clearpage
